\documentclass[11pt,a4paper]{article}
\usepackage[T1]{fontenc}
\usepackage[utf8]{inputenc}
\usepackage[brazilian]{babel}
\usepackage[default,type1,semibold]{inter}
\usepackage[a4paper,top=18mm,bottom=24mm,left=18mm,right=18mm,footskip=12mm]{geometry}
\usepackage[table]{xcolor}
\usepackage{amssymb,textcomp,array,longtable,enumitem,fancyhdr,lastpage,needspace,fontawesome5}
\usepackage[skins,breakable,raster]{tcolorbox}
\usepackage[hidelinks]{hyperref}
\usetikzlibrary{positioning,calc,babel}
\DeclareUnicodeCharacter{2192}{\ensuremath{\rightarrow}}
\DeclareUnicodeCharacter{2212}{\ensuremath{-}}
\DeclareUnicodeCharacter{2265}{\ensuremath{\geq}}
\DeclareUnicodeCharacter{2248}{\ensuremath{\approx}}
\DeclareUnicodeCharacter{03C3}{\ensuremath{\sigma}}

\definecolor{cor}{HTML}{2563EB}
\definecolor{suave}{HTML}{EAF0FE}
\definecolor{texto}{HTML}{0F172A}
\definecolor{cinza}{HTML}{5B6B84}
\definecolor{linha}{HTML}{D3DAE6}
\color{texto}
\setlength{\parindent}{0pt}
\setlength{\parskip}{4pt}
\linespread{1.1}
\setlength{\LTpre}{6pt}
\setlength{\LTpost}{10pt}
\renewcommand{\arraystretch}{1.45}
\arrayrulecolor{linha}
\setlist{itemsep=4pt,topsep=4pt,parsep=0pt}
\setlist[itemize]{label={\color{cor}\small\faCaretRight},leftmargin=1.4em}
% Passos numerados: número dentro de um círculo da cor da equipe.
\DeclareRobustCommand{\bolinha}[1]{\tikz[baseline=(n.base)]\node[circle,fill=cor,text=white,inner sep=1.6pt,minimum size=17pt,font=\footnotesize\bfseries](n){#1};}
\setlist[enumerate]{label=\bolinha{\arabic*},leftmargin=2.3em,labelsep=0.6em,itemsep=6pt}
% Checklist: quadrado grande para marcar à mão.
\DeclareRobustCommand{\caixinha}{\tikz[baseline=-0.6ex]\draw[cor,line width=1.2pt,rounded corners=2pt](0,-0.16)rectangle(0.42,0.26);}
\newlist{checklist}{itemize}{1}
\setlist[checklist]{label=\caixinha,leftmargin=2em,labelsep=0.6em,itemsep=7pt}

\pagestyle{fancy}
\fancyhf{}
\renewcommand{\headrulewidth}{0pt}
\fancyfoot[L]{\footnotesize\color{cor}\textbf{\faProjectDiagram\ \ Frente A}\color{cinza}\ \ ·\ \ Organização do Projeto · Guia do projeto}
\fancyfoot[R]{\footnotesize\color{cinza}\textbf{\thepage}\,/\,\pageref{LastPage}}
% Fita da cor da equipe na margem esquerda de todas as páginas.
\fancyhead[L]{\begin{tikzpicture}[remember picture,overlay]\fill[cor](current page.north west)rectangle([xshift=5pt]current page.south west);\end{tikzpicture}}

% Título de seção com ícone: nunca sozinho no pé da página.
\newcommand{\secao}[2]{\Needspace{10\baselineskip}\vspace{12pt}%
  \noindent\tikz[baseline=(i.base)]\node[circle,fill=cor,text=white,minimum size=26pt,inner sep=0pt](i){\small #1};%
  \hspace{8pt}{\Large\bfseries #2}\par\vspace{-4pt}{\color{cor!35}\rule{\linewidth}{1.2pt}}\par\vspace{2pt}}
\newcommand{\subsecao}[1]{\Needspace{8\baselineskip}\vspace{8pt}{\large\bfseries\color{cor}#1}\par}
\newcommand{\rotulo}[2]{\Needspace{6\baselineskip}\vspace{10pt}%
  \noindent{\color{cor}\normalsize #1}\hspace{6pt}{\footnotesize\bfseries\color{cor}\MakeUppercase{#2}}\par\vspace{2pt}}
\newtcolorbox{destaque}[1][]{enhanced,breakable,colback=suave,colframe=suave,borderline west={4pt}{0pt}{cor},arc=4pt,boxrule=0pt,left=12pt,right=12pt,top=8pt,bottom=8pt,before skip=8pt,after skip=10pt,#1}
\newtcolorbox{apostila}{enhanced,breakable,colback=white,colframe=cor,boxrule=1pt,arc=6pt,left=12pt,right=12pt,top=8pt,bottom=8pt,before skip=8pt,after skip=8pt,
  overlay={\node[fill=cor,text=white,rounded corners=3pt,font=\scriptsize\bfseries,inner sep=4pt,anchor=west] at ([xshift=10pt]frame.north west) {\faBookOpen\ DA APOSTILA};}}
\newtcolorbox{selo}{enhanced,left=8pt,right=8pt,top=8pt,bottom=8pt,before skip=0pt,after skip=0pt,
  frame code={\draw[cor,dashed,line width=1.4pt,rounded corners=8pt](frame.south west)rectangle(frame.north east);},
  interior code={\fill[suave,rounded corners=8pt](interior.south west)rectangle(interior.north east);}}
% Anotações: ocupa o resto da página da etapa (só se sobrar espaço útil).
\newcommand{\anotacoes}{\par\ifdim\dimexpr\pagegoal-\pagetotal\relax>3.5cm\ifdim\pagegoal<\maxdimen
\begin{tcolorbox}[enhanced,height fill,colback=white,colframe=cor!35,boxrule=0.8pt,arc=8pt,before skip=10pt,left=10pt,top=6pt,
  borderline={0.8pt}{0pt}{cor!35,dashed},frame hidden,
  underlay={\begin{scope}\clip(interior.south west)rectangle(interior.north east);
    \foreach \y in {1,...,40}{\draw[cor!18,line width=0.6pt]([yshift=-\y*0.75cm-0.5cm]interior.north west)--([yshift=-\y*0.75cm-0.5cm]interior.north east);}\end{scope}}]
{\footnotesize\bfseries\color{cor}\faStickyNote\ \ ANOTAÇÕES}
\end{tcolorbox}\fi\fi}
% Fim da etapa: se o bloco final não cabe, a sobra vira Anotações e o bloco vai para a próxima página.
\newif\ifjaanotou
\newsavebox{\blocofinal}
% Mede o bloco final já montado; se não couber no resto da página, a sobra vira Anotações.
\newcommand{\usarblocofinal}{\par\global\jaanotoufalse
  \ifdim\pagegoal<\maxdimen\ifdim\dimexpr\pagegoal-\pagetotal\relax<\dimexpr\ht\blocofinal+\dp\blocofinal+18pt\relax
    \anotacoes\global\jaanotoutrue\clearpage\fi\fi
  \vspace{10pt}\noindent\usebox{\blocofinal}\par}
\newtcolorbox{faixa}[1][]{enhanced,colback=cor,colframe=cor,coltext=white,arc=8pt,boxrule=0pt,left=14pt,right=14pt,top=10pt,bottom=10pt,#1}
\newtcbox{\chip}[1][]{enhanced,on line,colback=suave,colframe=suave,coltext=cor,arc=7pt,boxrule=0pt,left=5pt,right=5pt,top=2pt,bottom=2pt,fontupper=\footnotesize\bfseries,#1}
\newtcbox{\chipbranco}{enhanced,on line,colback=white,colframe=white,coltext=cor,arc=7pt,boxrule=0pt,left=6pt,right=6pt,top=2pt,bottom=2pt,fontupper=\footnotesize\bfseries}
\begin{document}
\thispagestyle{fancy}
\begin{tikzpicture}[remember picture,overlay]
\fill[cor,opacity=0.10] ([xshift=-1.5cm,yshift=-1cm]current page.north east) circle (5.2cm);
\fill[cor,opacity=0.07] ([xshift=-5.5cm,yshift=-4.5cm]current page.north east) circle (2.4cm);
\fill[cor,opacity=0.06] ([xshift=2cm,yshift=3cm]current page.south west) circle (4cm);
\end{tikzpicture}
\vspace*{2mm}
\noindent\begin{tikzpicture}
\node[circle,fill=cor,text=white,minimum size=78pt,font=\fontsize{40}{40}\selectfont\bfseries] (L) {A};
\node[right=14pt of L.east,anchor=west,align=left,text width=12.2cm] {
  {\footnotesize\bfseries\color{cor}PROJETO ESTADIA\ \ ·\ \ EQUIPE AZUL\ \ ·\ \ FRENTE A}\\[4pt]
  {\fontsize{23}{27}\selectfont\bfseries Organização do Projeto}\\[6pt]
  \chip{\faCompass\ \ Guia do projeto}\ \chip{\faProjectDiagram\ \ Aulas 03, 06 e 08 da apostila}
};
\end{tikzpicture}
\vspace{10pt}

\begin{faixa}
{\large\faRocket}\hspace{8pt}{\large\bfseries Nossa missão}\par\vspace{2pt}
{\normalsize Organizar o trabalho do grupo: o que vamos fazer, quem faz, até quando, e juntar tudo no final.}
\end{faixa}

\vspace{6pt}
\begin{tcolorbox}[enhanced,colback=white,colframe=cor!40,boxrule=1pt,arc=8pt,left=10pt,right=10pt,top=8pt,bottom=10pt,
  title={\faUserFriends\ \ Nossa equipe},fonttitle=\bfseries\small,coltitle=cor,colbacktitle=suave,attach boxed title to top left={xshift=10pt,yshift=-8pt},boxed title style={arc=4pt,boxrule=0pt}]
\vspace{6pt}
\begin{minipage}[t]{0.31\linewidth}\centering
\tikz\node[circle,draw=cor,line width=1.4pt,fill=white,minimum size=34pt,text=cor]{\large \faCompass};\par\vspace{4pt}
{\scriptsize\bfseries\color{cor}\MakeUppercase{Coordenação}}\par\vspace{14pt}
{\color{cor!60}\rule{0.9\linewidth}{0.8pt}}\par{\scriptsize\color{cinza}nome}
\end{minipage}\hfill
\begin{minipage}[t]{0.31\linewidth}\centering
\tikz\node[circle,draw=cor,line width=1.4pt,fill=white,minimum size=34pt,text=cor]{\large \faClipboard};\par\vspace{4pt}
{\scriptsize\bfseries\color{cor}\MakeUppercase{Documentação}}\par\vspace{14pt}
{\color{cor!60}\rule{0.9\linewidth}{0.8pt}}\par{\scriptsize\color{cinza}nome}
\end{minipage}\hfill
\begin{minipage}[t]{0.31\linewidth}\centering
\tikz\node[circle,draw=cor,line width=1.4pt,fill=white,minimum size=34pt,text=cor]{\large \faLaptopCode};\par\vspace{4pt}
{\scriptsize\bfseries\color{cor}\MakeUppercase{Dev de apoio}}\par\vspace{14pt}
{\color{cor!60}\rule{0.9\linewidth}{0.8pt}}\par{\scriptsize\color{cinza}nome}
\end{minipage}
\end{tcolorbox}
\begin{destaque}[before skip=4pt]
{\color{cor}\Large\faCrosshairs}\hspace{8pt}{\large\bfseries O seu foco}\par\vspace{3pt}
Vocês organizam o grupo: o que precisa ser feito, quem faz e até quando. Não precisam entender de programação.
\end{destaque}
\secao{\faSitemap}{Como usar este guia}
{\large Este guia tem uma parte geral e um \textbf{manual para cada função}. Leia a parte geral e depois vá direto ao manual da sua função: ele tem tudo de que você precisa para fazer as suas tarefas do Passo a passo.}\par\vspace{6pt}
\Needspace{7\baselineskip}
\begin{tcbraster}[raster columns=3,raster equal height=rows,raster column skip=6pt,raster row skip=6pt,raster before skip=4pt,raster after skip=8pt]
\begin{tcolorbox}[enhanced,colback=white,colframe=cor!25,boxrule=0.8pt,arc=6pt,left=8pt,right=8pt,top=6pt,bottom=6pt,colback=suave!40]
\centering{\color{cor}\LARGE \faCompass}\par\vspace{3pt}{\bfseries Coordenação}\par\vspace{2pt}{\footnotesize\color{cinza}Conduz as reuniões e acompanha se o grupo está no prazo.}\par\vspace{4pt}\chip{A2}\ \chip{A3}\par\vspace{4pt}{\footnotesize\bfseries\color{cor}\faArrowRight\ manual na página \pageref{manual-a1}}
\end{tcolorbox}
\begin{tcolorbox}[enhanced,colback=white,colframe=cor!25,boxrule=0.8pt,arc=6pt,left=8pt,right=8pt,top=6pt,bottom=6pt,colback=suave!40]
\centering{\color{cor}\LARGE \faClipboard}\par\vspace{3pt}{\bfseries Documentação}\par\vspace{2pt}{\footnotesize\color{cinza}Registra o projeto no papel: Canvas, ata e o trabalho final.}\par\vspace{4pt}\chip{A1}\ \chip{A4}\par\vspace{4pt}{\footnotesize\bfseries\color{cor}\faArrowRight\ manual na página \pageref{manual-a2}}
\end{tcolorbox}
\begin{tcolorbox}[enhanced,colback=white,colframe=cor!25,boxrule=0.8pt,arc=6pt,left=8pt,right=8pt,top=6pt,bottom=6pt,colback=suave!40]
\centering{\color{cor}\LARGE \faLaptopCode}\par\vspace{3pt}{\bfseries Dev de apoio}\par\vspace{2pt}{\footnotesize\color{cinza}Mostra o sistema e confirma as informações técnicas.}\par\vspace{4pt}\chip{apoio às tarefas}\par\vspace{4pt}{\footnotesize\bfseries\color{cor}\faArrowRight\ manual na página \pageref{manual-a-dev}}
\end{tcolorbox}
\end{tcbraster}
\secao{\faHotel}{O projeto em 1 minuto}
{\large O Estadia é um sistema para administrar um hotel de 50 quartos. Pelo computador ou pelo celular, a equipe do hotel faz reservas, registra a entrada (check-in) e a saída (check-out) dos hóspedes, cobra o que eles consumiram, fecha o caixa do dia e vê relatórios.}\par
\Needspace{7\baselineskip}
\begin{tcbraster}[raster columns=2,raster equal height=rows,raster column skip=6pt,raster row skip=6pt,raster before skip=4pt,raster after skip=8pt]
\begin{tcolorbox}[enhanced,colback=white,colframe=cor!25,boxrule=0.8pt,arc=6pt,left=8pt,right=8pt,top=6pt,bottom=6pt,colback=suave!50]
{\color{cor}\Large \faExclamationTriangle}\hspace{6pt}{\bfseries O problema}\par\vspace{3pt}{\small Muitos hotéis pequenos ainda controlam tudo em caderno, planilha e WhatsApp. Aí acontecem erros: vender o mesmo quarto para duas pessoas, cobrar errado, o caixa não bater e a dona não saber quanto o hotel realmente ganhou no mês.}
\end{tcolorbox}
\begin{tcolorbox}[enhanced,colback=white,colframe=cor!25,boxrule=0.8pt,arc=6pt,left=8pt,right=8pt,top=6pt,bottom=6pt,colback=suave!50]
{\color{cor}\Large \faLightbulb}\hspace{6pt}{\bfseries A solução}\par\vspace{3pt}{\small Um sistema único que vários funcionários usam ao mesmo tempo. Ele não deixa vender o mesmo quarto duas vezes, guarda quem fez cada ação e mostra relatórios prontos para a dona decidir preços e equipe.}
\end{tcolorbox}
\end{tcbraster}
\subsecao{\faAngleDoubleRight\ O que vocês precisam saber}
\begin{destaque}\begin{itemize}
  \item O sistema já está pronto (versão 1.0). O trabalho de vocês é organizar a parte de gestão do projeto.
  \item São 4 frentes: A (vocês), B (produto e tarefas), C (metas e indicadores) e D (pessoas e treinamento).
  \item Cada frente entrega seus arquivos no site da divisão de funções. É de lá que vocês juntam tudo no final.
  \item Se uma frente atrasar, avisem cedo: o SPI serve exatamente para isso.
\end{itemize}
\end{destaque}
\subsecao{\faAngleDoubleRight\ Riscos que vale citar no Canvas}
\begin{destaque}\begin{itemize}
  \item Faltar tempo de alguém do grupo (provas, trabalho).
  \item A dona do hotel demorar para responder perguntas.
  \item Os funcionários do hotel terem dificuldade de usar o sistema.
  \item A internet do hotel cair (o sistema continua funcionando, porque roda na rede local).
\end{itemize}
\end{destaque}
\clearpage
\begin{faixa}[before skip=0pt]
\begin{tabular}{@{}m{46pt}>{\raggedright\arraybackslash}m{0.76\linewidth}@{}}
\tikz\node[circle,fill=white,text=cor,minimum size=42pt,inner sep=0pt]{\LARGE \faCompass}; &
{\footnotesize\bfseries MANUAL DA FUNÇÃO · FRENTE A}\label{manual-a1}\par\vspace{2pt}
{\fontsize{20}{24}\selectfont\bfseries Coordenação}\par\vspace{3pt}
{\small Conduz as reuniões e acompanha se o grupo está no prazo.}
\end{tabular}
\end{faixa}
\rotulo{\faTasks}{Suas tarefas}
\chip{A2}\ Reunião de início (kick-off)\hspace{10pt}\chip{A3}\ Cronograma e acompanhamento\par
{\footnotesize\color{cinza}Você também ajuda nas outras tarefas da frente. Quem faz cada uma está no Passo a passo.}\par
\Needspace{12\baselineskip}\vspace{10pt}
\begin{tcolorbox}[enhanced,colback=cor,colframe=cor,coltext=white,arc=6pt,boxrule=0pt,left=10pt,right=10pt,top=6pt,bottom=6pt,after skip=6pt]
{\Large \faFlag}\hspace{10pt}{\footnotesize\bfseries TAREFA A2}\hspace{8pt}{\large\bfseries Reunião de início (kick-off)}\hfill{\footnotesize\faClock\ cerca de 1h (30 min de reunião)}
\end{tcolorbox}
\rotulo{\faBrain}{O que você precisa saber}
\begin{itemize}
  \item O kick-off não é para descobrir o projeto: é para todos saírem sabendo o que fazer.
  \item Toda decisão anotada precisa de um responsável e uma data.
\end{itemize}
\Needspace{7\baselineskip}
\begin{tcbraster}[raster columns=2,raster equal height=rows,raster column skip=6pt,raster row skip=6pt,raster before skip=4pt,raster after skip=8pt]
\begin{tcolorbox}[enhanced,colback=white,colframe=cor!25,boxrule=0.8pt,arc=6pt,left=8pt,right=8pt,top=6pt,bottom=6pt,colback=suave!50]
{\footnotesize\bfseries\color{cor}\MakeUppercase{Dicas}}\par\vspace{3pt}{\small {\color{cor}\faLightbulb}\ \ Coordenação fala, Documentação anota: não tentem fazer as duas coisas ao mesmo tempo.}
\end{tcolorbox}
\begin{tcolorbox}[enhanced,colback=white,colframe=cor!25,boxrule=0.8pt,arc=6pt,left=8pt,right=8pt,top=6pt,bottom=6pt,colback=suave!50]
{\footnotesize\bfseries\color{cor}\MakeUppercase{Cuidado com}}\par\vspace{3pt}{\small {\color{cor}\faExclamationTriangle}\ \ Terminar sem datas. ``Vamos ver depois'' não é combinado.}
\end{tcolorbox}
\end{tcbraster}
\subsecao{\faBoxOpen\ Roteiro da reunião (30 minutos)}
\begingroup\small
\begin{longtable}{>{\raggedright\arraybackslash}p{0.105\linewidth}>{\raggedright\arraybackslash}p{0.823\linewidth}}
\rowcolor{cor} \color{white}\textbf{\footnotesize Minutos} & \color{white}\textbf{\footnotesize O que fazer} \\
\endfirsthead
\rowcolor{cor} \color{white}\textbf{\footnotesize Minutos} & \color{white}\textbf{\footnotesize O que fazer} \\
\endhead
\rowcolor{white}\textbf{0–5} & Boas-vindas e objetivo do trabalho. \\ \hline
\rowcolor{suave!55}\textbf{5–10} & O que é o Estadia (use ``O projeto em 1 minuto'' deste guia). \\ \hline
\rowcolor{white}\textbf{10–20} & Mostrar o site: as 4 frentes, as tarefas e como enviar entregas. \\ \hline
\rowcolor{suave!55}\textbf{20–25} & Cronograma das 4 semanas. \\ \hline
\rowcolor{white}\textbf{25–30} & Combinados e próximos passos (ler a ata em voz alta). \\ \hline
\end{longtable}
\endgroup
\subsecao{\faBoxOpen\ Mensagem pronta para convidar o grupo}
\begin{destaque}{\color{cor}\faQuoteLeft}\ \ Oi, pessoal! Reunião de início do trabalho do Estadia: [dia] às [hora], 30 minutos, [link]. Antes, escolham sua vaga no site da divisão de funções: [link do site].\end{destaque}
\Needspace{12\baselineskip}\vspace{10pt}
\begin{tcolorbox}[enhanced,colback=cor,colframe=cor,coltext=white,arc=6pt,boxrule=0pt,left=10pt,right=10pt,top=6pt,bottom=6pt,after skip=6pt]
{\Large \faCalendar}\hspace{10pt}{\footnotesize\bfseries TAREFA A3}\hspace{8pt}{\large\bfseries Cronograma e acompanhamento}\hfill{\footnotesize\faClock\ cerca de 1h + 15 min por semana}
\end{tcolorbox}
\rotulo{\faBrain}{O que você precisa saber}
\begin{itemize}
  \item Linha de base = a primeira versão do cronograma, guardada sem mudar.
  \item SPI = tarefas prontas ÷ tarefas que deveriam estar prontas. 1 = no prazo; abaixo de 1 = atrasado; acima de 1 = adiantado.
\end{itemize}
\Needspace{7\baselineskip}
\begin{tcbraster}[raster columns=2,raster equal height=rows,raster column skip=6pt,raster row skip=6pt,raster before skip=4pt,raster after skip=8pt]
\begin{tcolorbox}[enhanced,colback=white,colframe=cor!25,boxrule=0.8pt,arc=6pt,left=8pt,right=8pt,top=6pt,bottom=6pt,colback=suave!50]
{\footnotesize\bfseries\color{cor}\MakeUppercase{Dicas}}\par\vspace{3pt}{\small {\color{cor}\faLightbulb}\ \ No Google Sheets, a coluna SPI pode ser a fórmula =C2/B2.\par\vspace{3pt}{\color{cor}\faLightbulb}\ \ Atualizem toda sexta-feira, no mesmo horário.}
\end{tcolorbox}
\begin{tcolorbox}[enhanced,colback=white,colframe=cor!25,boxrule=0.8pt,arc=6pt,left=8pt,right=8pt,top=6pt,bottom=6pt,colback=suave!50]
{\footnotesize\bfseries\color{cor}\MakeUppercase{Cuidado com}}\par\vspace{3pt}{\small {\color{cor}\faExclamationTriangle}\ \ Mudar as datas da linha de base quando algo atrasa. Ela fica igual; quem muda é a ``data real''.}
\end{tcolorbox}
\end{tcbraster}
\subsecao{\faBoxOpen\ Todas as entregas do grupo (para o cronograma)}
\begingroup\small
\begin{longtable}{>{\raggedright\arraybackslash}p{0.105\linewidth}>{\raggedright\arraybackslash}p{0.823\linewidth}}
\rowcolor{cor} \color{white}\textbf{\footnotesize Semana} & \color{white}\textbf{\footnotesize Entregas} \\
\endfirsthead
\rowcolor{cor} \color{white}\textbf{\footnotesize Semana} & \color{white}\textbf{\footnotesize Entregas} \\
\endhead
\rowcolor{white}\textbf{1} & A1 Canvas · A2 Ata do kick-off · B1 Backlog · B2 Kanban · C1 OKR · D1 Conversa com a dona \\ \hline
\rowcolor{suave!55}\textbf{2} & A3 Cronograma · B3 Sprint · C2 Indicadores · D2 Manuais \\ \hline
\rowcolor{white}\textbf{3} & B4 Reuniões · C3 Painel · D3 Treinamento \\ \hline
\rowcolor{suave!55}\textbf{4} & D4 Quem decide o quê · A4 Trabalho final \\ \hline
\end{longtable}
\endgroup
\subsecao{\faBoxOpen\ Três exemplos de SPI}
\begingroup\small
\begin{longtable}{>{\raggedright\arraybackslash}p{0.241\linewidth}>{\raggedright\arraybackslash}p{0.143\linewidth}>{\raggedright\arraybackslash}p{0.077\linewidth}>{\raggedright\arraybackslash}p{0.395\linewidth}}
\rowcolor{cor} \color{white}\textbf{\footnotesize Deveriam estar prontas} & \color{white}\textbf{\footnotesize Estão prontas} & \color{white}\textbf{\footnotesize SPI} & \color{white}\textbf{\footnotesize Significa} \\
\endfirsthead
\rowcolor{cor} \color{white}\textbf{\footnotesize Deveriam estar prontas} & \color{white}\textbf{\footnotesize Estão prontas} & \color{white}\textbf{\footnotesize SPI} & \color{white}\textbf{\footnotesize Significa} \\
\endhead
\rowcolor{white}\textbf{6} & 6 & 1,00 & No prazo \\ \hline
\rowcolor{suave!55}\textbf{8} & 6 & 0,75 & Atrasado: conversar com quem atrasou \\ \hline
\rowcolor{white}\textbf{5} & 6 & 1,20 & Adiantado \\ \hline
\end{longtable}
\endgroup
\clearpage
\begin{faixa}[before skip=0pt]
\begin{tabular}{@{}m{46pt}>{\raggedright\arraybackslash}m{0.76\linewidth}@{}}
\tikz\node[circle,fill=white,text=cor,minimum size=42pt,inner sep=0pt]{\LARGE \faClipboard}; &
{\footnotesize\bfseries MANUAL DA FUNÇÃO · FRENTE A}\label{manual-a2}\par\vspace{2pt}
{\fontsize{20}{24}\selectfont\bfseries Documentação}\par\vspace{3pt}
{\small Registra o projeto no papel: Canvas, ata e o trabalho final.}
\end{tabular}
\end{faixa}
\rotulo{\faTasks}{Suas tarefas}
\chip{A1}\ Canvas do projeto\hspace{10pt}\chip{A4}\ Fechamento e apresentação\par
{\footnotesize\color{cinza}Você também ajuda nas outras tarefas da frente. Quem faz cada uma está no Passo a passo.}\par
\Needspace{12\baselineskip}\vspace{10pt}
\begin{tcolorbox}[enhanced,colback=cor,colframe=cor,coltext=white,arc=6pt,boxrule=0pt,left=10pt,right=10pt,top=6pt,bottom=6pt,after skip=6pt]
{\Large \faThLarge}\hspace{10pt}{\footnotesize\bfseries TAREFA A1}\hspace{8pt}{\large\bfseries Canvas do projeto}\hfill{\footnotesize\faClock\ cerca de 1h30}
\end{tcolorbox}
\rotulo{\faBrain}{O que você precisa saber}
\begin{itemize}
  \item Canvas = o projeto inteiro numa página. Cada bloco responde uma pergunta simples.
  \item Uma frase por bloco basta. Se precisar de parágrafo, a ideia ainda não está clara.
\end{itemize}
\Needspace{7\baselineskip}
\begin{tcbraster}[raster columns=2,raster equal height=rows,raster column skip=6pt,raster row skip=6pt,raster before skip=4pt,raster after skip=8pt]
\begin{tcolorbox}[enhanced,colback=white,colframe=cor!25,boxrule=0.8pt,arc=6pt,left=8pt,right=8pt,top=6pt,bottom=6pt,colback=suave!50]
{\footnotesize\bfseries\color{cor}\MakeUppercase{Dicas}}\par\vspace{3pt}{\small {\color{cor}\faLightbulb}\ \ No Canva, procurem o modelo ``Project Model Canvas''.\par\vspace{3pt}{\color{cor}\faLightbulb}\ \ Usem a cor da Frente A (azul) no quadro.}
\end{tcolorbox}
\begin{tcolorbox}[enhanced,colback=white,colframe=cor!25,boxrule=0.8pt,arc=6pt,left=8pt,right=8pt,top=6pt,bottom=6pt,colback=suave!50]
{\footnotesize\bfseries\color{cor}\MakeUppercase{Cuidado com}}\par\vspace{3pt}{\small {\color{cor}\faExclamationTriangle}\ \ Objetivo sem verbo (``sistema do hotel''). Comece com um verbo: ``Implantar…'', ``Organizar…''.}
\end{tcolorbox}
\end{tcbraster}
\subsecao{\faBoxOpen\ Informações do projeto para usar nos blocos}
\begingroup\small
\begin{longtable}{>{\raggedright\arraybackslash}p{0.166\linewidth}>{\raggedright\arraybackslash}p{0.762\linewidth}}
\rowcolor{cor} \color{white}\textbf{\footnotesize Bloco} & \color{white}\textbf{\footnotesize Informação que ajuda} \\
\endfirsthead
\rowcolor{cor} \color{white}\textbf{\footnotesize Bloco} & \color{white}\textbf{\footnotesize Informação que ajuda} \\
\endhead
\rowcolor{white}\textbf{Objetivo} & Um sistema para o hotel controlar quartos, reservas e caixa sem erros e sem depender de internet. \\ \hline
\rowcolor{suave!55}\textbf{Benefícios} & Fim do overbooking, caixa que bate, relatórios para a dona decidir preço. \\ \hline
\rowcolor{white}\textbf{Interessados} & Dona do hotel, recepcionistas, governança, gerência, hóspedes, contador, o grupo e o professor. \\ \hline
\rowcolor{suave!55}\textbf{Entregas} & O sistema (pronto) + os documentos das 4 frentes + apresentação. \\ \hline
\rowcolor{white}\textbf{Riscos} & Falta de tempo do grupo, dona demorar a responder, funcionários com dificuldade. \\ \hline
\rowcolor{suave!55}\textbf{Prazo} & 4 semanas (ver o cronograma do site). \\ \hline
\rowcolor{white}\textbf{Equipe} & 12 pessoas: 4 devs e 4 duplas de gestão. \\ \hline
\rowcolor{suave!55}\textbf{Como medir} & As metas da Frente C (ex.: check-in em 3 minutos). \\ \hline
\end{longtable}
\endgroup
\Needspace{12\baselineskip}\vspace{10pt}
\begin{tcolorbox}[enhanced,colback=cor,colframe=cor,coltext=white,arc=6pt,boxrule=0pt,left=10pt,right=10pt,top=6pt,bottom=6pt,after skip=6pt]
{\Large \faFlagCheckered}\hspace{10pt}{\footnotesize\bfseries TAREFA A4}\hspace{8pt}{\large\bfseries Fechamento e apresentação}\hfill{\footnotesize\faClock\ cerca de 2h}
\end{tcolorbox}
\rotulo{\faBrain}{O que você precisa saber}
\begin{itemize}
  \item O trabalho final junta as entregas das 4 frentes, na ordem A, B, C e D.
\end{itemize}
\Needspace{7\baselineskip}
\begin{tcbraster}[raster columns=2,raster equal height=rows,raster column skip=6pt,raster row skip=6pt,raster before skip=4pt,raster after skip=8pt]
\begin{tcolorbox}[enhanced,colback=white,colframe=cor!25,boxrule=0.8pt,arc=6pt,left=8pt,right=8pt,top=6pt,bottom=6pt,colback=suave!50]
{\footnotesize\bfseries\color{cor}\MakeUppercase{Dicas}}\par\vspace{3pt}{\small {\color{cor}\faLightbulb}\ \ Baixem as entregas pelo site: a mais recente fica no topo de cada frente.\par\vspace{3pt}{\color{cor}\faLightbulb}\ \ Ensaiem com cronômetro: cerca de 2 minutos por frente.}
\end{tcolorbox}
\begin{tcolorbox}[enhanced,colback=white,colframe=cor!25,boxrule=0.8pt,arc=6pt,left=8pt,right=8pt,top=6pt,bottom=6pt,colback=suave!50]
{\footnotesize\bfseries\color{cor}\MakeUppercase{Cuidado com}}\par\vspace{3pt}{\small {\color{cor}\faExclamationTriangle}\ \ Colar as entregas sem revisar: confiram se todas as tabelas estão preenchidas.}
\end{tcolorbox}
\end{tcbraster}
\subsecao{\faBoxOpen\ Estrutura do documento final}
\begin{destaque}\begin{enumerate}
  \item Capa e sumário.
  \item Introdução: o que é o Estadia (1 parágrafo).
  \item Frente A: Canvas, ata e cronograma com SPI.
  \item Frente B: backlog, Kanban, sprint e reuniões.
  \item Frente C: OKR, indicadores e painel.
  \item Frente D: conversa com a dona, manuais, treinamento e governança.
  \item Lições aprendidas e conclusão.
\end{enumerate}
\end{destaque}
\subsecao{\faBoxOpen\ Slides sugeridos (10 a 12)}
\begin{destaque}\begin{itemize}
  \item 1 abertura
  \item 1 o que é o Estadia
  \item 2 por frente (8 no total)
  \item 1 lições aprendidas
  \item 1 encerramento
\end{itemize}
\end{destaque}
\Needspace{22\baselineskip}\vspace{14pt}
\begin{faixa}[before skip=0pt]
\begin{tabular}{@{}m{46pt}>{\raggedright\arraybackslash}m{0.76\linewidth}@{}}
\tikz\node[circle,fill=white,text=cor,minimum size=42pt,inner sep=0pt]{\LARGE \faLaptopCode}; &
{\footnotesize\bfseries MANUAL DA FUNÇÃO · FRENTE A}\label{manual-a-dev}\par\vspace{2pt}
{\fontsize{20}{24}\selectfont\bfseries Dev de apoio}\par\vspace{3pt}
{\small Mostra o sistema e confirma as informações técnicas.}
\end{tabular}
\end{faixa}
\secao{\faTools}{O que você entrega em cada tarefa}
\begingroup\small
\begin{longtable}{>{\raggedright\arraybackslash}p{0.105\linewidth}>{\raggedright\arraybackslash}p{0.823\linewidth}}
\rowcolor{cor} \color{white}\textbf{\footnotesize Tarefa} & \color{white}\textbf{\footnotesize O que fazer} \\
\endfirsthead
\rowcolor{cor} \color{white}\textbf{\footnotesize Tarefa} & \color{white}\textbf{\footnotesize O que fazer} \\
\endhead
\rowcolor{white}\textbf{A1} & Explicar o sistema em 5 minutos e conferir ``Entregas'' e ``Equipe'' do Canvas. \\ \hline
\rowcolor{suave!55}\textbf{A2} & Participar do kick-off e explicar o que já está pronto no sistema. \\ \hline
\rowcolor{white}\textbf{A3} & Avisar se alguma tarefa técnica pode atrasar. \\ \hline
\rowcolor{suave!55}\textbf{A4} & Revisar os slides que falam do sistema. \\ \hline
\end{longtable}
\endgroup
\rotulo{\faHandsHelping}{Sempre}
\begin{itemize}
  \item Mostrar o sistema funcionando e criar usuários de teste.
  \item Conferir o que a frente escreveu sobre o sistema.
  \item Passar números e prints da versão de demonstração.
  \item Explicar o que for técnico em palavras simples.
\end{itemize}
\begin{destaque}{\color{cor}\faDesktop}\ \ Peça ao seu dev de apoio o endereço da versão de demonstração e um usuário de teste. Ela tem dados inventados: nunca use dados reais de hóspedes nos trabalhos ou prints.\end{destaque}
\clearpage
\secao{\faQuestionCircle}{Onde tirar dúvidas}
\begingroup\small
\begin{longtable}{>{\raggedright\arraybackslash}p{0.342\linewidth}>{\raggedright\arraybackslash}p{0.586\linewidth}}
\rowcolor{cor} \color{white}\textbf{\footnotesize Preciso de…} & \color{white}\textbf{\footnotesize Onde encontrar} \\
\endfirsthead
\rowcolor{cor} \color{white}\textbf{\footnotesize Preciso de…} & \color{white}\textbf{\footnotesize Onde encontrar} \\
\endhead
\rowcolor{white}\textbf{Ver o sistema funcionando} & Peça ao dev de apoio o link da versão de demonstração \\ \hline
\rowcolor{suave!55}\textbf{Tirar dúvida sobre o sistema} & Dev de apoio da sua frente \\ \hline
\rowcolor{white}\textbf{Conceitos da disciplina} & Apostila ``Gestão Estratégica Guiada a Dados'' \\ \hline
\rowcolor{suave!55}\textbf{O que a dona do hotel respondeu} & Frente D (questionário da reunião) \\ \hline
\rowcolor{white}\textbf{Entregar seu trabalho} & Site da divisão de funções → Entregas da sua frente \\ \hline
\end{longtable}
\endgroup
\secao{\faBook}{Glossário}
\Needspace{27\baselineskip}
\begin{tcbraster}[raster columns=2,raster equal height=rows,raster column skip=6pt,raster row skip=6pt,raster before skip=4pt,raster after skip=8pt]
\begin{tcolorbox}[enhanced,colback=white,colframe=cor!25,boxrule=0.8pt,arc=6pt,left=8pt,right=8pt,top=6pt,bottom=6pt]
{\bfseries\color{cor}Diária}\par\vspace{1pt}{\footnotesize O preço de uma noite no hotel.}
\end{tcolorbox}
\begin{tcolorbox}[enhanced,colback=white,colframe=cor!25,boxrule=0.8pt,arc=6pt,left=8pt,right=8pt,top=6pt,bottom=6pt]
{\bfseries\color{cor}Check-in / Check-out}\par\vspace{1pt}{\footnotesize Registro da entrada e da saída do hóspede.}
\end{tcolorbox}
\begin{tcolorbox}[enhanced,colback=white,colframe=cor!25,boxrule=0.8pt,arc=6pt,left=8pt,right=8pt,top=6pt,bottom=6pt]
{\bfseries\color{cor}Taxa de ocupação}\par\vspace{1pt}{\footnotesize Quantos \% dos quartos foram ocupados. Ex.: 35 de 50 = 70\%.}
\end{tcolorbox}
\begin{tcolorbox}[enhanced,colback=white,colframe=cor!25,boxrule=0.8pt,arc=6pt,left=8pt,right=8pt,top=6pt,bottom=6pt]
{\bfseries\color{cor}Diária média (ADR)}\par\vspace{1pt}{\footnotesize Quanto o hotel recebeu, em média, por quarto vendido.}
\end{tcolorbox}
\begin{tcolorbox}[enhanced,colback=white,colframe=cor!25,boxrule=0.8pt,arc=6pt,left=8pt,right=8pt,top=6pt,bottom=6pt]
{\bfseries\color{cor}RevPAR}\par\vspace{1pt}{\footnotesize Receita das diárias dividida por todos os quartos (vendidos ou não).}
\end{tcolorbox}
\begin{tcolorbox}[enhanced,colback=white,colframe=cor!25,boxrule=0.8pt,arc=6pt,left=8pt,right=8pt,top=6pt,bottom=6pt]
{\bfseries\color{cor}No-show}\par\vspace{1pt}{\footnotesize Quem reservou e não apareceu.}
\end{tcolorbox}
\begin{tcolorbox}[enhanced,colback=white,colframe=cor!25,boxrule=0.8pt,arc=6pt,left=8pt,right=8pt,top=6pt,bottom=6pt]
{\bfseries\color{cor}Governança}\par\vspace{1pt}{\footnotesize Equipe de limpeza e arrumação dos quartos.}
\end{tcolorbox}
\begin{tcolorbox}[enhanced,colback=white,colframe=cor!25,boxrule=0.8pt,arc=6pt,left=8pt,right=8pt,top=6pt,bottom=6pt]
{\bfseries\color{cor}Backlog}\par\vspace{1pt}{\footnotesize Lista de tudo que o sistema precisa ter, em ordem de importância.}
\end{tcolorbox}
\begin{tcolorbox}[enhanced,colback=white,colframe=cor!25,boxrule=0.8pt,arc=6pt,left=8pt,right=8pt,top=6pt,bottom=6pt]
{\bfseries\color{cor}Sprint}\par\vspace{1pt}{\footnotesize Um ciclo curto de trabalho (aqui, 1 semana) com um objetivo.}
\end{tcolorbox}
\begin{tcolorbox}[enhanced,colback=white,colframe=cor!25,boxrule=0.8pt,arc=6pt,left=8pt,right=8pt,top=6pt,bottom=6pt]
{\bfseries\color{cor}KPI}\par\vspace{1pt}{\footnotesize Indicador importante, que ajuda a tomar uma decisão.}
\end{tcolorbox}
\end{tcbraster}
\end{document}
